% Ad Familiares 4.8 --- headnote
% ad-familiares-04-08, mid-August 46 BC, Rome

\textit{Cicero to M. Marcellus, written at Rome in the latter
half of August 46 BC --- Perseus: \textit{Romae in m.~Sext.~a.~708
(46)}. This is the short follow-up note that falls between the
two substantial letters of 13 August (4.7 and 4.9) and the next
substantial letter that survives (4.9 if read in its other
guise; 4.10 in November). The position of this letter within
the cluster is documentary rather than rhetorical: Cicero is
keeping the channel to Mytilene open while the case is taking
shape at Rome that will end with the speech \textit{Pro Marcello}
in September.}

\textit{The texture is still Cicero's, but the argument is
trimmed almost to a slogan. The familiar Ciceronian formula ---
\textit{si sit aliqua res publica \dots sin autem nulla sit} ---
is bent here into a single dilemma with no third term: either
the republic exists, in which case Marcellus must be present
in it; or it does not, in which case Rome is at least the most
fitting place to go into exile, since exile from one tyranny
into another part of the same tyranny accomplishes nothing.
The text at the close of section~1 is corrupt in the manuscripts
(\textit{omnia debere †tua causa\dots}); the translation here
follows the standard editorial sense.}
